% Template de Relatório para Trabalho de Fases Paralelas
% Compilar com: pdflatex report_template.tex

\documentclass[10pt,twocolumn]{article}

% Pacotes
\usepackage[utf8]{inputenc}
\usepackage[portuguese]{babel}
\usepackage[margin=2cm]{geometry}
\usepackage{graphicx}
\usepackage{amsmath}
\usepackage{booktabs}
\usepackage{listings}
\usepackage{xcolor}
\usepackage{hyperref}

% Configurações
\setlength{\columnsep}{0.6cm}
\setlength{\parindent}{0pt}
\setlength{\parskip}{0.3em}

% Título
\title{\textbf{Bubble Sort Paralelo com Modelo de Fases Paralelas usando MPI}}
\author{
    Nome do Aluno 1 \\
    \texttt{email1@edu.pucrs.br}
    \and
    Nome do Aluno 2 \\
    \texttt{email2@edu.pucrs.br}
}
\date{\today}

\begin{document}

\maketitle

\section{Introdução}

Este trabalho implementa o algoritmo Bubble Sort seguindo o modelo de \textbf{fases paralelas} usando a biblioteca MPI (Message Passing Interface). O objetivo é analisar o desempenho desta abordagem e compará-la com o modelo de divisão e conquista para o mesmo problema.

\section{Metodologia}

\subsection{Modelo de Fases Paralelas}

O modelo de fases paralelas funciona da seguinte forma:

\begin{enumerate}
    \item \textbf{Distribuição:} Cada processo é responsável por $\frac{1}{np}$ do vetor global
    \item \textbf{Fase 1 - Ordenação Local:} Cada processo ordena seu subvetor usando Bubble Sort
    \item \textbf{Fase 2 - Verificação:} Cada processo verifica se está ordenado em relação ao vizinho (comunicação point-to-point + broadcast)
    \item \textbf{Fase 3 - Troca de Valores:} Se não estiver ordenado, processos trocam valores com vizinhos para convergir
    \item Repete fases 1-3 até convergência global
\end{enumerate}

\subsection{Ambiente de Execução}

Os experimentos foram executados no cluster \textbf{Grad} com as seguintes configurações:

\begin{itemize}
    \item \textbf{Nós:} 2 nós exclusivos
    \item \textbf{Processos:} 16 cores (32 com Hyper-Threading)
    \item \textbf{Tamanho do vetor:} 1.000.000 elementos
    \item \textbf{Dados:} Ordem decrescente (pior caso)
    \item \textbf{Implementação MPI:} OpenMPI
\end{itemize}

\section{Resultados}

\subsection{Tempos de Execução}

A Tabela~\ref{tab:results} apresenta os tempos de execução para diferentes configurações:

\begin{table}[h]
\centering
\small
\begin{tabular}{@{}ccccc@{}}
\toprule
\textbf{Procs} & \textbf{Tempo (s)} & \textbf{Speedup} & \textbf{Efic. (\%)} & \textbf{Iter.} \\
\midrule
1 (seq)  & XX.XXX & 1.00  & 100.0 & - \\
4        & XX.XXX & X.XX  & XX.X  & XX \\
8        & XX.XXX & X.XX  & XX.X  & XX \\
16       & XX.XXX & X.XX  & XX.X  & XX \\
32       & XX.XXX & X.XX  & XX.X  & XX \\
\bottomrule
\end{tabular}
\caption{Resultados experimentais (1M elementos)}
\label{tab:results}
\end{table}

\textbf{Observações:}
\begin{itemize}
    \item Speedup calculado como: $S_p = \frac{T_{seq}}{T_{par}}$
    \item Eficiência calculada como: $E_p = \frac{S_p}{p} \times 100\%$
    \item Número de iterações indica convergência do algoritmo
\end{itemize}

\subsection{Análise de Desempenho}

\textbf{Speedup e Escalabilidade:}

[PREENCHER COM ANÁLISE DOS RESULTADOS]

\textit{Exemplo:} O speedup observado foi de X.XX com 16 processos, indicando uma eficiência de XX\%. A escalabilidade foi [boa/moderada/limitada] devido a [overhead de comunicação/sincronização/etc].

\textbf{Número de Iterações:}

O algoritmo convergiu em aproximadamente XX iterações para todas as configurações, indicando que [análise sobre comportamento de convergência].

\section{Comparação com Divisão e Conquista}

A Tabela~\ref{tab:comparison} compara os dois modelos:

\begin{table}[h]
\centering
\small
\begin{tabular}{@{}lccc@{}}
\toprule
\textbf{Modelo} & \textbf{Tempo (s)} & \textbf{Speedup} & \textbf{Efic.} \\
\midrule
Divisão/Conquista & XX.XXX & X.XX & XX\% \\
Fases Paralelas   & XX.XXX & X.XX & XX\% \\
\bottomrule
\end{tabular}
\caption{Comparação entre modelos (16 processos)}
\label{tab:comparison}
\end{table}

\textbf{Análise Comparativa:}

[PREENCHER COM COMPARAÇÃO DETALHADA]

\textit{Exemplo:} O modelo de fases paralelas apresentou [melhor/pior] desempenho comparado ao modelo de divisão e conquista. Isso se deve a [características específicas de cada modelo: padrão de comunicação, balanceamento de carga, overhead, etc].

\textbf{Vantagens do Modelo de Fases Paralelas:}
\begin{itemize}
    \item Distribuição uniforme de dados (1/np por processo)
    \item Padrão de comunicação regular e previsível
    \item Boa localidade de dados após convergência
\end{itemize}

\textbf{Desvantagens:}
\begin{itemize}
    \item Múltiplas iterações até convergência
    \item Overhead de verificação global a cada iteração
    \item Trocas de mensagens frequentes entre vizinhos
\end{itemize}

\section{Conclusões}

[PREENCHER COM CONCLUSÕES]

Este trabalho demonstrou a implementação e análise do algoritmo Bubble Sort usando o modelo de fases paralelas com MPI. Os resultados mostraram que [principais conclusões sobre desempenho, escalabilidade, comparação com outros modelos].

Trabalhos futuros poderiam investigar [otimizações, outras aplicações do modelo, etc].

\section*{Código Fonte}

O código fonte completo está disponível em:

\texttt{https://github.com/seu-usuario/mpi-training}

\vspace{0.5em}

\textbf{Arquivos principais:}
\begin{itemize}
    \item \texttt{seq.c} - Versão sequencial
    \item \texttt{mpi\_bubblesort\_phases.c} - Fases paralelas
    \item \texttt{run\_phases\_experiments.sh} - Scripts de testes
\end{itemize}

\end{document}
